\chapter{Metodologia}
\label{c:Metodologia}
Després d'haver fet la recerca previa, i d'haver decidit fer una xarxa neuronal de regressió que compleixi tots els requesits que esperavem, ara cal pensar amb quina metodologia podrem aconseguir el nostre objectiu. Per tant vol dir que hem de fer un altre recerca en el que utilizarem les millors eines per fer el Treball de Recerca de manera ràpida i eficient.\\


Primer ens em fet la pregunta de  sobre quins àmbits hem de recorre, una vegada feta la pregunta ja buscarem les millors eines o mètode que ens ajudarà en cadascún dels àmbits. A continuació farem unes presentacions dels àmbits que vam tractar i les diferents solucions que vam donar i el definitiu.\\

Per començar tractem  l'apartat \ref{sec:4.1}, que consisteix en com vam estructurar la via de comunicació entre nosaltres i el tutor. A l'apartat \ref{sec:4.2}, tracta de quin editor de text vam escollir, que es un punt crucial per el treball ja que les eines que ens aporta l'editor en si mateix ens pot estalviar molt de temps i convertirlo visualment més accessible. A continuació hi es l'apartat \ref{sec:4.3}, on t'explica la plataforma que hem escollit per l'entorn col·laboratiu. Finalment per poder crear una xarxa neuronal hem d'utilizar un llenguatge de programació i això s'explicarà a l'apartat \ref{sec:4.4}.




 \section{Comunicació} \label{sec:4.1}
 Com que som dos persones fent el treball amb un tutor, ens és fonamental tenir una bona via comunicativa. Per tant, hem acollit diferents maneres de comunicació. Ara bé no em pogut registrar tots els registres que hem tingut però la gran majoria s'ha pogut conservar.\\ \\
 D'aquí cap avall explicarem les diferents maneres de comunicació utilizades en l'elaboracio d'aquest treball.\\
    \begin{enumerate}
     \item \textbf{Correu electrònic:} L'eina que vam utilizar a distància és el servei de correu electrònic Gmail. Vam fer ús d'aquest aparell per el fet de que és simple i eficient, i té una gran memòria d'emmagatzament per conservar tots els correus que ens hem fet.
     Les converses fetes per el correu electrònic ha sigut la principal via de comunicació que em establert, per la facilitat i l'eficiència que ens aporta. Per tant, el tema ensenyament i dubtes simples o feiem tot per gmail.
      \item \textbf{Git:} El Git~\cite{git} és un sistema de repositoris que permenten el treball col$\cdot$laboratiu i està especialment dissenyat per als projectes informàtics. En el nostre TR hem creat un projecte (TR-RJ-2025~\cite{TR-RJ-2025}) a
      la plataforma GitHub~\cite{GitHub} que ofereix una plataforma que a desenvolupar projectes de codi obert.
      %%% En un paràgraf 3 citacions, això és documentar!!
      \item \textbf{Fulls de càlculs:} Hem realitzat diverses reunions amb el tutor, i la via que vam utilitzar per configurar l'horari de les quedades és a travès d'un full de càlcul. En aquestes quedades ens em organitzat el treball i ens em quedat amb els diferents mètodes que utilizariem per el trebal, el tutor ens arregla qualsevol tipus de problema, i la instalació dels programes necessaris.
    \end{enumerate}


 \section{Editor de text i processador de text: }\label{sec:4.2}
 Per aconseguir una bona presentació de TR, haviem d'escollir o un editor de text o un processador. Tant l'editor com el processador són per crear documents, pero cadascún te les seves funcions.\\ \\
 L'editor de text, és tal com el seu nom, serveix per editar els documents. És molt simple i això genera inconvenients, com per exemple no poder utilizar formats avançat o una personificació. Però, les avantatge d'aquest és que pots tenir la capacitat de navegar rapidament de fitxer en fitxer, facilita la codificació i té compabilitat amb tots els tipus de llenguatge. És l'ideal per programadors.\\ \\
 Un processador de text són tipus de programes més elaborats i més complexos. Com que són programes amb tanta complexibilitat, et dona accès a formats avançat, utilitzat per a documents professionals com currículums, llibres i d'altres més, també et donen tot tipus d'opcions col·laboratives com disseny de pàgina, correció gràfica i molt més.\\ \\
 Per a la part teòrica del nostre projecte sobre la Intel·ligència Artificial, vam seleccionar el format PDF degut a:
    \begin{enumerate}
      \item \textbf{Connexió segura i estable:} Permet l'accés sense la dependència d'una connexió a Internet constant.
      \item \textbf{Baixos requisits de recursos:} Pot obrir-se en gairebé qualsevol dispositiu sense necessitat de programes específiques.
      \item \textbf{Portabilitat:} Manté el format independentment del sistema operatiu utilitzat.
      \item \textbf{Multifuncionalitat:} Admet text, fórmules matemàtiques, imatges i codi de programació integrat.
    \end{enumerate}
  L'estalvi de temps va ser un factor decisiu, ja que evita problemes de compatibilitat entre fórmules matemàtiques, imatges i el format del document.
  En la nostra recerca d'eines que satisfacin tots els requisits, vam identificar dos tipus principals de processadors: Una que es deia WYSIWYG (What You See Is What You Get) i l'altre el WYSIWYM (What You See Is What You Mean). Ara donarem una explicació de les advantatges i desavantatges que tenen:

 \begin{enumerate}
  \item \textbf{What You See Is What You Get(WYSIWYG):} WYSIWYG, es refereix a un tipus de processador que permet als usuaris veure en temps real el resultat final del document o disseny mentre editen, sense la necessitat de coneixer el codi o llenguatges de marcatges. Aquests tipus d'editors poden ser tant com el microsoft word o Google drive.
  \textbf{Característiques}:
    \begin{enumerate}
     \item \textbf{Interfície visual:} Els canvis de format (negretes, taules, imatges) es fan mitjançant eines gràfiques, no mitjançant codi.
     \item \textbf{Generació automàtica de codi:} El programa crea el codi subjacent (HTML, CSS, etc.) sense que l'usuari hi intervingui directament.
     \item \textbf{Ús generalitzat:} S'utilitza en àmbits com l'edició web, el disseny gràfic i la publicació digital.
    \end{enumerate}

  \textbf{Avantatges}:
    \begin{enumerate}
     \item \textbf{Edició visual intuïtiva:} Permet ajustar el format directament (arrossegar imatges, canviar fonts amb clics).
     \item Ideal per a usuaris sense coneixements tècnics.
     \item \textbf{Resultat immediat:} No cal compilar, els canvis es veuen al moment.
     \item \textbf{Eines integrades:} Funcions com correcció ortogràfica, taules gràfiques, o opcions de disseny accessibles des del menú
     \item \textbf{Millor per a maquetació complexa:} Documents amb molts elements gràfics (pòsters, fulls informatius)
    \end{enumerate}

  \textbf{Desavantatges}:
    \begin{enumerate}
      \item \textbf{Poca precisió en contingut tècnic:} Fórmules matemàtiques o codis es poden deformar-se a l'hora de canviar el format.
     \item \textbf{Errors en documents llargs:} El format manual pot provocar errors (salt de pàgina, numeració desorganitzada).
     \item \textbf{Dependència del programa:} Si s'obre en un altre software, el disseny pot variar.
     \item \textbf{Control de versions complicat:} Difícil treballar en equip sense conflictes de format.
    \end{enumerate}
 \item \textbf{What You See Is What You Mean(WYSIWYM):}\label{subsec:4.2.2} WSYSIWYM, es refereix a un tipus de processador que es centra en l'estructura del contigut, no en la seva aparença visual immediata, l'usuari marca el text segons la seva funció (títol, secció, cita i d'altres més) i el format final s'aplica mitjançant un full d'estil com css o Latex.

 \textbf{Avantatges:}
    \begin{enumerate}
     \item \textbf{Precisió en elements tècnics:} Fórmules matemàtiques, codi o referències es gestionen amb sintaxi clara.
     \item \textbf{Consistència automàtica:} L'estil s'aplica globalment.
     \item \textbf{Lleuger i portable:} Els fitxers són de text pla.
     \item \textbf{Ideal per a treball en entorns col·laboratius:} Es pot actualizar amb git i fusionar sense haver de canviar el format.
    \end{enumerate}
 \textbf{Desavantatges:}
    \begin{enumerate}
     \item \textbf{Corba d'aprenentatge:} Cal aprendre una sintaxi específica.
     \item \textbf{Previsualització no immediata:} En LaTeX, cal compilar; en Markdown, cal un renderitzador.
     \item \textbf{Limitacions en disseny gràfic:} Personalització avançada (ex: posicionament exacte d'imatges) requereix codi addicional.
    \end{enumerate}
      \textbf{Exemples:}
         \begin{enumerate}

          \item \textbf{\LaTeX:} \LaTeX, és un tipus de WYSIWYM, on s'enfoca en l'estructura lògica, no la visual, tal com vam mencionar en l'apartat \ref{subsec:4.2.2}. LaTeX és un sistema de composició tipogràfica d'alta qualitat, especialment dissenyat per a la creació de documents acadèmics, tècnics i científics. A diferència dels processadors de text tradicionals com Microsoft Word, LaTeX es basa en codi per estructurar el document, separant el contingut de la seva presentació visual. Va ser desenvolupat per Leslie Lamport en 1984 com un conjunt de macros per al sistema TeX de Donald Knuth. A més, \LaTeX{} és \nameref{subsec:programari lliure}, distribuït sota la LaTeX Project Public License (LPPL), cosa que permet utilitzar-lo, modificar-lo i redistribuir-lo lliurement. La raó principal de la seva creació va ser la impotència dels processadors de text de l'època per compilar fórmules matemàtiques avançades i gestionar documents complexos amb precisió tipogràfica. És per aquesta raó la qual el vam escollir com a processador per satisfer els requeriments que demanem. Per més informació o dubte podeu consultar el manual de latex: \href{https://manualdelatex.com/}{Manual de LaTeX}

          \item \textbf{Kile:} Kile es un editor especific del \LaTeX  que funciona en els sistemes operatius tant Linux, Windows com Apple Macintosh. L'editor va ser desenvolupat per la comunitat  KDE\cite{KDE} en que ofereix diverses eines avannçades per a l'edicio de LaTeX. No és l'únic editor que existeix, però és el que més ens convé per el fet de que sigui creat per documents acadèmics i científics.\\
          \textbf{Avantatges:}
              \begin{enumerate}
              \item \textbf{Entorn integrat complet:} Kile ofereix tot ho necessari per treballar amb LaTeX en una sola interfície, incloent editor, compilador, visualitzador i gestor de referències
              \item \textbf{Altament configurable: } Permet afegir botons a la barra d'eines per executar scripts personalitzats i automatitzar tasques complexes.
              \item \textbf{Feedback WYSWYG:} Et dona una sensació quasi WYSWYG, gracies als visualitzadors que els té integrades en que pots veure els canvis en temps reals.
              \item \textbf{Eines avançades d'edició:} Té autocompletador de comandes, accés ràpid a símbols matemàtics, navegacions ràpides per seccions.
              \end{enumerate}

          \textbf{Desavantatges:}
              \begin{enumerate}
               \item \textbf{Nivell d'apranentatge:} Costa molt d'aprendre l'interfície si no estas familialitzat amb KDE.
               \item \textbf{Configuració inicial:} Requereix molts tipus de configuracions externs per tal d'arrenca.
              \end{enumerate}
          \begin{comment}
          Ara mostrarem una taula de comparació respecte els altres editors:
         LyX és més adequat per a principiants gràcies a la seva interfície WYSIWYG. No obstant, Kile es més flexible per als usuaris avançats i te més control en els codis i funcions. \\
         %Taula 1
         \begin{table}[h!]
          \begin{tabular}{|l|l|l|}
         \hline
          \LaTeX & \textbf{Kile} & \textbf{LyX} \\ \hline
         \textbf{Facilitat d'ús} &  & X \\ \hline
          \textbf{Control i Flexibilitat} & X & \\ \hline
          \textbf{Visualització en temps real} & & X \\ \hline
         \end{tabular}
         \end{table}

         TeXmaker i TeXstudio comparteixen moltes similituds amb Kile en termes de funcionalitat i d'interfície. \\

         %Taula 2
         \begin{table}[h!]
          \begin{tabular}{|l|l|l|}
         \hline
          \LaTeX & \textbf{Kile} & \textbf{TeXmaker i TeXstudio} \\ \hline
          \textbf{similitud de funcionalitat} & - & - \\ \hline
          \textbf{Integració KDE} & X & \\ \hline
          \textbf{Disponibilitat dels sistemes operatius} & X &  \\ \hline
          \end{tabular}
         \end{table}

         Gummi és un editor molt més simple que Kile, amb previsualització en temps real però amb moltes menys funcionalitats. És bo per a documents simples però per a projectes complexos no. \\
         %Taula 3
         \begin{table}[h!]
          \begin{tabular}{|l|l|l|}
         \hline
          \LaTeX & \textbf{Kile} & \textbf{Gummi} \\ \hline
          \textbf{complexibilitat:}  & X &  \\ \hline
          \textbf{Previsualització} &  & X \\ \hline
          \textbf{funcionalitat} & X &  \\ \hline
          \end{tabular}
         \end{table}

         Per més informació en qüestió de comparació podeu fer una consulta a les següents pagines: \href{https://osluca.uca.es/noticia/cinco-editores-de-latex-libres/}{5 editors de LaTeX} \href{https://latex.org/forum/viewtopic.php?t=208}{Analisí detallada de Kile i d'altres editors}
         Fonts: \href{https://iloo.wordpress.com/2010/10/20/kile-otro-editor-latex/}{Experiencia Personal de l'us Kile} \href{https://www.kdeblog.com/editor-de-latex-para-kde-kile.html}{Descripcio de Característiques Kile}
         \end{comment}
         \begin{comment}
         \item \textbf{KDE:} El nom "KDE"\ originalment sí que era un acrònim de "K Desktop Environment" (des de la seva creació el 1996), però a partir del 2009, la comunitat va decidir deixar de considerar-lo un acrònim i utilitzar-lo simplement com un nom propi per representar:
             \begin{enumerate}
               \item L'entorn d'escriptori.
               \item La comunitat global que desenvolupa programari lliure.
               \item Tots els projectes relacionats.
             \end{enumerate}
          \end{comment}

         \end{enumerate}

\end{enumerate}




\begin{table}[h!]
\begin{tabular}{|l|l|l|}
\hline
\textbf{Aspecte} & \textbf{WYSIWYG} & \textbf{WYSIWYM} \\ \hline
\textbf{Enfocament} & Aparença visual immediata & Estructura semàntica del contingut \\ \hline
\textbf{Exemples} & Microsoft Word, WordPress & LaTeX, LyX \\ \hline
\textbf{Usuaris} & No tècnics, dissenyadors &  Acadèmics, desenvolupadors tècnics \\ \hline
\textbf{Control} &  Limitada (codi generat automàtic) & Alt (definició manual de l'estructura) \\ \hline
\end{tabular}
\caption{Comparativa entre WYSIWYG i WYSIWYM}
\end{table}



font(WYSIWYG): (\href{https://es.wikipedia.org/wiki/WYSIWYG}{Wikipedia} \href{https://es.wikipedia.org/wiki/WYSIWYG}{Arimetrics} \href{https://www.wix.com/encyclopedia/definition/wysiwyg}{Wix}).
font(WYSIWYM): (\href{https://en.wikipedia.org/wiki/WYSIWYM}{Wikipedia} \href{https://en.ryte.com/wiki/WYSIWYG/}{Ryte WIKI}).

\begin{comment}
% l'entorn figure et numera la figura i et permet posar una nota al peu de la figura, però costa controlar la posició
\begin{figure}[h!]
    \centering
    \includegraphics[width=0.5\textwidth]{./figures/latex.png}
    \caption{Logo del Latex}
\end{figure}%
¨
% si ho fas d'aqesta manera
\begin{center}
    \includegraphics[width=0.75\textwidth]{./figures/latex.pdf}
\end{center}
\end{comment}

\subsubsection{Programari lliure}\label{subsec:programari lliure}

El programari (o software en anglès) és el conjunt de programes informàtics, procediments i documentació que permeten a un ordinador realitzar tasques específiques. És la part interna d'una computadora ,a difèrencia del hardware que és l'extern, el software és part física.\\

Un programari es considera Free software quan respecta les quatres llibertats fonamentals definides per la Free Software Foundation (FSF):
\begin{enumerate}
 \item  \textbf{Llibertat 0:} Executar el programa amb qualsevol propòsit.
 \item  \textbf{Llibertat 1:} Estudiar i modificar el codi font (requereix accés al codi)
 \item  \textbf{Llibertat 2:} Redistribuir còpies per ajudar altres usuaris.
 \item  \textbf{Llibertat 3:} Millorar el programa i compartir les modificacions.
\end{enumerate}

Font:\href{https://www.fsf.org/}{Free Software Foundation}
\begin{comment}
\subsubsection{L'entorn d'escriptori}
Un entorn d'escriptori en sistemes operatius és la interfície gràfica que permet als usuaris interactuar amb l'ordinador. Inclou elements com finestres, icones, panells, fons de pantalla i eines de gestió d'arxius. A continuació, mostrarem els principals entorns que n'hi ha:
\begin{enumerate}
 \item KDE Plasma
 \begin{itemize}
  \item \textbf{Característiques: } Altament personalitzable, amb temes globals, widgets i configuració avançada.
  \item \textbf{Recomanat per:} Usuaris que volen control total sobre l'aparença i el flux de treball.
 \end{itemize}

 \item GNOME
 \begin{itemize}
  \item \textbf{Característiques:} Disseny minimalista i enfocat en la productivitat. Compatible amb extensions per ampliar funcionalitats.
  \item \textbf{Recomanat per: }  Usuaris que prefereixen una experiència neta i senzilla.
 \end{itemize}


 \item XFCE
\begin{itemize}
  \item \textbf{Característiques:} Lleuger i ràpid, ideal per a maquinari antic o limitat.
  \item \textbf{Recomanat per: } Usuaris que busquen eficiència sense sacrificar funcionalitats bàsiques .
 \end{itemize}

 \item Cinnamon
\begin{itemize}
  \item \textbf{Característiques:} Interfície tradicional similar a Windows, amb personalització mitjana.
  \item \textbf{Recomanat per: } Usuaris que provenen de Windows i busquen una transició suau.
 \end{itemize}

\end{enumerate}

Fonts:\href{https://dev.to/xploitcore/kde-vs-gnome-vs-others-choosing-the-best-linux-desktop-environment-in-2025-ab5}{Dev} \href{https://planet.communia.org/index.php/en/node/77}{Planet}
\end{comment}


\section{Entorn col·laboratiu}\label{sec:4.3}
A l'hora de crear un projecte en equip, pot arribar a ser complicat treballar amb 2 o més membres del grup a la vegada en els seus respectius dispositius si no s'escull una eina col·laborativa adequada, sobretot si s'està treballant amb molts fitxers. Pot haver problemes alhora de intercanviar o que sol pugui estar un membre traballant en un fitxer fa que sigui un absolutament caos.
\begin{enumerate}
 \item \textbf{Google Drive}
 Google Drive és una plataforma de Google que ofereix eines gratuïtes com un processador de text (Docs) i un entorn col·laboratiu. Tot i ser senzill i eficient, té algunes limitacions importants que hem de considerar:
\begin{itemize}
 \item \textbf{Poc control en el procés col·laboratiu:} És difícil gestionar canvis simultanis o conflicts entre versions.
 \item \textbf{Limitacions en el control de versions:} No permet controlar detalladament els canvis que es fan i registrar.
 \item \textbf{Falta de privadesa:} Google té accés als teus arxius, cosa que pot ser un problema per a dades sensibles.
\end{itemize}
Aquestes limitacions fan que Google Drive no sigui l'opció ideal per el nostre TR la qual cosa ens fara escollir l'altre alternativa que teniem, el Git+GitHub+Vim.

\item \textbf{Git+GitHub+Vim}

Git és un sistema de control de versions que ofereix:
\begin{itemize}
 \item Historial dels canvis.
 \item Branques per desenvolupar diferents funcionalitat sense afectar la versió principal.
 \item Control total dels canvis
 \item Control i revisió de codi
\end{itemize}
El principal desavantatge que té el Git és que, per si sol, no està optimitzat per a la col·laboració en equips grans. Pero aquí entra en acció el GitHub, que complementa Git amb:
\begin{itemize}
\item \textbf{Repositoris}
 \item \textbf{Git pull,commit,push:} Dona la possibilitat de poder discutir les versions abans de pujarlo a la original.
 \item Solucions als conflictes de versions amb l'ajuda del Vim~\cite{Vim}.
 \end{itemize}

\textbf{Avantatges:}

\begin{itemize}
 \item \textbf{Copia de seguretat:} Cada usuari té una copia de seguretat del servidor principal, en el cas d'errors o corrupció, sempre el pots reemplaçar-lo per el servidor principal.
 \item \textbf{Garantia de dades:} El model de dades que utilitza Git garantitza la integritat ciptográfica de cada bit del teu projecte. Cada arxiu es recupera mitjançant la seva suma de comprobació quan es torna a desgarregar-se. Es imposible obtenir res de Git que no sigui els bits exactes que s'ha integrat.
 \item \textbf{Àrea de preparació:} Una cosa que diferència a Git d'altres eines es que es posible preparar alguns dels teus arxius i confirmar-los sense confirmar tots els altres canvis que s'han fet en els altres arxius, això et permet preparar nomès algunes parts d'un arxiu modificat i es realitza amb el commandament git commit ``nom de l'arxiu''. Git també et dona la opció d'ignorar aquesta característica i confirmar tots els canvis que s'han fet, afegint un ``-a'' al final del commandament.
 \item \textbf{Eficiència de Git:} Git no depèn d'internet per la majoria de les operacions, totes les accions principals (veure l'historial, fer commits) es fan en localment. Git està programat en llenguatge C, això permet a Git realitzar les tasques amb molta rapidesa encara que es treballi amb molts fichers.
\end{itemize}
Més informació en: \href{https://git-scm.com/about/branching-and-merging}{Pagina oficial del Git}

\textbf{Desavantatges:}
\begin{itemize}



\item \textbf{Corba d'aprenentatge pronunciada}: Git té un model conceptual complex (working directory, staging area, repositori local/remot) que requereix temps per assimilar. Operacions bàsiques com merge/resolucions de conflictes poden resultar difícils per a principiants.

      \begin{enumerate}
        \item \textbf{Working directory:}
         \begin{itemize}
            \item La carpeta local del projecte on treballes i fas canvis als fitxers.
            \item Modifiques, elimines o crees nous arxius
            \item Git analiza aquest aquest directori per fer la comparació dels canvis respecte l'ultima versió
          \end{itemize}

        \item \textbf{Staging Area:}
         \begin{itemize}
            \item És el cache del Git
            \item Una zona intermedia on inclous els canvis propers que vols fer
          \end{itemize}

        \item \textbf{Repositori Local:}
        \begin{itemize}
            \item És la base de dades completa de Git que es troba a la teva màquina (dins de la carpeta .git)
            \item Conté tot l'historial de commits, branques, tags del Git
            \item Els canvis del Staging Area, es guarden permanent aqui
          \end{itemize}
        \item \textbf{Repositori Remot:}
        \begin{itemize}
            \item És una copia del repositori en un servidor (GitHub, GitLab, Bitbucket)
            \item Serveix per compartir codis i funciona com a entorn col·laboratiu
            \item És utilizen comandes git push(Pujar canvis) i git pull(Baixar canvis)
          \end{itemize}
       \end{enumerate}


\item \textbf{Problemes amb fitxers molt grans}: Git no està optimitzat per a fitxers binaries de gran mida (vídeos, imatges, bases de dades). El rendiment es degrada significativament i el repositori pot inflar-se massa.

\item \textbf{Gestió complicada de credencials}: Git no gestiona de forma òptima les credencials.

\item \textbf{Problemes de rendiment en repositoris molt grans}: En repositoris amb milers de commits o arxius, certes operacions (com git status o git checkout) poden tornar-se lentes.

\item \textbf{Comandes}: Els comandes reben nom contraintuïtiu io la sintaxi varia segons l'operació.
\end{itemize}
Fonts:\href{https://ieeexplore.ieee.org/document/8259410}{Comparacions entre control de versions} \href{https://stevebennett.me/2012/02/24/10-things-i-hate-about-git/}{bloc de critica d'un desenvolupador}
\end{enumerate}

\begin{comment}
\subsubsection{Vim}\label{subsubsec:Vim}
Vim (Vi IMproved) és un editor de text avançat, lliure i de codi obert, basat en l'editor clàssic vi. És conegut per ser multiplataforma i també la seva eficiència, personalització i ús sense ratoli, sent popular entre programadors i usuaris tècnics. Aquest editor de text consta de 4 modes amb diferents funcionalitats:
\begin{itemize}
 \item \textbf{Mode Estandard:} Per navegar i executar ordres
 \item \textbf{Mode Insert:} Per escriure text
 \item \textbf{Mode Visual:} Per seleccionar bloc de text
 \item \textbf{Mode Command} Per executar comandes complexes
\end{itemize}
Fonts:\href{https://www.vim.org/}{Pagina oficial del Vim} \href{https://vimhelp.org/}{Documentacio del Vim}
\end{comment}
\section{Llenguatge de programació}\label{sec:4.4}

El llenguatge de programació que vam decidir per fer la part pràctia del treball va ser Python. Python és un llenguatge de programació que va ser creat en 1989 per Guido van Rossum , un informàtic neerlandès. Aquest llenguatge és conegut per la seva simple sintaxi, comparada amb altres llenguatges com JavaScript o C++ que són més complicats, python facilita l'aprenentatge pels usuaris que començen a apendre a programar desde 0.

\subsection{Python}
La raó principal per la que vam escollir aquest llenguatgera per les avantatges seguents:

\textbf{Avantatges}
\begin{itemize}
 \item \textbf{No cal compilar:} No és necesari compilar en Python, un cop que tens el codi, s'executa directament.
 \item \textbf{Llibreries} Python té una gran quantitat de llibreries molt útils com ``TensorFlow'' o numpy, aquestes llibreries estalvien moltes línies de codi a l'hora de programar una xarxa neuronal.
 \item \textbf{Identificació d'errors} Encara que altres llenguatges també poden identificar-te els erros, Python t'ho explica de manera més clara i precisa, estalviant-te el temps d'anar línea per línea buscant l'error.
 \item \textbf{Sintaxi:}La seva sintaxi clara i llegible que està explicada a l'apartat \ref{4.4.2} és ho que fa que aquest llenguatge sigui fàcil d'apendre quan un comença a programar desde 0.
\end{itemize}

Malgrat totes aquestes avantatges, Python tambè té unes porques desaventatges que hem de tenir present.

\textbf{Desavantatges:}
\begin{itemize}
 \item \textbf{Limitacions de rendiment:} La naturalesa interpretada de Python significa que és més lent que altres llenguatges. La causa d'això és que Python es tradueix a codi màquina línea per línea durant l'execució del codi, cosa que afegeix un temps de processament.
 \item \textbf{Consum de memoria:} Gràcies a la bona fexibilitat d'aquest llenguate fa que l'ordinador requereixi un major consum de memòria, un consum excesiu de la memoria pot causar relentizació en el programa i bloquejar-lo. Per tant es recomendable tenir un ordinador o portatil decent que pugui executar correctament Python.
\end{itemize}

\subsection{Sintaxi}
Acabem d'explicar que la sintaxi de Python és simple, però què diferencia aquesta sintaxi d'altres llenguatges? La principal raó és que el seu format de codi és visualment ordenat i utilitza paraules clau de l'anglès. A diferència d'altres llenguatges, no utilitza claudàtors per determinar blocs, com per exemple les sentències ``if, else'':


\begin{table}[h!]
\centering
\begin{tabular}{|p{7cm}|p{7cm}|}
\hline
\textbf{Condicions en C} & \textbf{Condicions en Python} \\ \hline

\begin{minipage}[t]{\linewidth}
\begin{lstlisting}[language=C]
#include <stdio.h>

int main() {
    int numero = 5;

    if (numero > 0) {
        printf("És >0.\n");
    } else {
        printf("És <0.\n");
    }
    return 0;
}
\end{lstlisting}
\end{minipage}
&
\begin{minipage}[t]{\linewidth}
\begin{lstlisting}[language=Python]
numero = 5

if numero > 0:
    print("És positiu.")
else:
    print("És negatiu.")
\end{lstlisting}
\end{minipage}
\\ \hline
\end{tabular}
\caption{Comparació de condicions en C i Python}
\end{table}




En aquesta taula podem veurem la diferència de sintaxi que té cada llenguatge, en aquest cas podem apreciar que C és mès llarg que python i que requereis més commandaments i línies de codi. A continuació mostrarem una taula resumida d'aquestes diferències.
\\

\textbf{Resum de diferències:}
\begin{table}[h!]
\centering
\begin{tabular}{|l|l|l|}
\hline
\textbf{Concepte} & \textbf{C} & \textbf{Python} \\
\hline
Terminador de línia & \verb|;| obligatori & No és obligatori \verb|;| \\
\hline
Agrupació de blocs & \{ \} & Espai en blanc \\
\hline
Condició & \verb|if (condició)| & \verb|if condició:| \\
\hline
Estructura & Requereix el script \verb|main()| & Script directe \\
\hline
Librerías & \verb|#include <stdio.h>| & No requereix per \verb|print()| \\
\hline
\end{tabular}
\caption{Resum de les diferències}
\end{table}

\subsection{Editor de Text}
Per començar a programar, és molt important escollir bé l'editor de text. Per raons tècniques, ara farem una llistat de quines opcions teníem i quina vam escollir finalment. \\

\textbf{VS Code(Visual Studio Code):}
VS Code és un editor de codi gratuït i de codi obert desenvolupat per Microsoft, i continua sent una de les eines més populars en 2025. Es caracteritza per ser multiplataforma, lleuger, altament personalitzable i suporta tots els llenguatges de programació gràcies a les seves extensions.
Caracteristiques:
\begin{itemize}
 \item És caracteritza per tenir un control de versions integrat del Git.
 \item Inclou un depurador de codis pas per pas.
 \item Interfície intuïtiva modernitzat.
\end{itemize}
\textbf{La nostra experiència utilitzant VS Code:}
Vam començar a utilitzar VS Code perquè ja el dominàvem abans del Treball de Recerca (TR), cosa que representava un gran avantatge inicial. Tanmateix, ens vam trobar amb diversos problemes a l'hora d'executar certs codis. Aquests inconvenients,vam haber de buscar una alternativa,el Eclipse. \\

\textbf{Eclipse:}
Eclipse és un entorn integrat de desenvolupament (IDE) de codi obert, escrit principalment en Java, però dissenyat per treballar amb múltiples llenguatges de programació. originalment fou creada per international Business Machines (IBM). Actualement cau en mans del Fundació Eclipse. Eclipse te les mateixes característiques que el de VS Code,no obstant, té una funcionalitat de més que es el enfocament en JAVA i ofereix diverses extensions per ell i té una gestio de fitxers grans millors que el de VS  Code. \\ \\
\textbf{La nostra experiència utilitzant Eclipse:}
Nosaltres el vam escollir per la seva funcionalitat i també per el sugeriment pel nostre tutor, consta d'un rendiment i estabilitatg superior a la de VS Code, de fet te uns extensions més solides que la de VS Code, però al final l'hem hagut de deixar per el fet de que vam decidir treballar amb Python i no ens convenia massa el seu punt d'enfoc ja que pot funcionar molt bé amb arxius molts grans de JAVA, però d'altres no. Aquesta vegada hem estat fent una gran recerca d'ús d'editor de text i finalment vam trobar el Pycharm. \\ \\
\textbf{Pycharm:}
PyCharm, desenvolupat per JetBrains, és l'entorn integrat de desenvolupament (IDE) més popular per a Python en l'actualitat. Empreses com Twitter, Facebook, Amazon i Pinterest el fan servir per al desenvolupament de les seves aplicacions en Python, gràcies a les seves eines avançades i la seva gran adaptabilitat.\\
Característiques:
\begin{itemize}
 \item Diverses extensions per Python
 \item Integració amb frameworks populars
 \item Suport per a bases de dades
 \item Depuració avançada
\end{itemize}
\textbf{La nostra experiència utilitzant Pycharm:}
Després de fer una recerca molt amplia, finalment vam optar Pycharm. Després d'estar utilitzant molt de temps ens vam donar compte la excepcionalitat que conte en el enfocament en Python, un dels aspectes que més ens va sorprendre va ser la seva capacitat per estalviar-nos temps significatiu en el desenvolupament. La integració amb els principals frameworks de Python ens va permetre accelerar moltíssim tasques que abans requerien una configuració manual extensa.  Totes les avantatges que ens aporta al treball amb Python  ens va convecer en utilitzar-ho com a l'editor de text definitiu que treballarem al llarg del treball.\\
Fonts:\href{https://www.arsys.es/blog/que-es-visual-studio-code-y-cuales-son-sus-ventajas}{VS Code i les seves avantatges} \href{https://www.eclipse.org/}{Fundació Eclipse} \href{https://datascientest.com/es/pycharm}{Tot l'ho que has de saber del Pycharm}\href{https://www.jetbrains.com/pycharm/whatsnew/#}{Totes les novetats de Pycharm (Pagina oficial)}

\section{Sistema Operatiu}
Els programes i programaris d'alt nivell moltes vegades no són compatibles amb tots tipus de sistemes operatius, o que no tenen el mateix rendiment un respecte a l'altre. Per tant això ens porta a fer una recerca.\\ \\
Un sistema operatiu és:\\ \\
{\color{gray}``Para entender qué es un sistema operativo, podemos pensar en él como en un administrador central que gestiona todos los recursos del sistema. Actúa como intermediario entre el hardware de un dispositivo informático (como una computadora, un teléfono inteligente o una tableta) y las aplicaciones de software que se ejecutan en él. Controla las operaciones del hardware y facilita la ejecución de múltiples tareas y procesos a través de una interfaz sencilla. ``}\\ \\
 Tal com està escrita en la Universitat Europea~\cite{UniversitatEuropea}, llavors es classifiquen en diverses sistemes operatius elaborada per diferents companya i les més destacades són les següents: \nameref{subsec:Microsoft Windows} \nameref{subsec:Linux} \nameref{subsec:Mac OS}

\subsection{Windows}

\label{subsec:Microsoft Windows}
Windows és un sistema operatiu desenvolupat per Microsoft a principis dels anys 80. La seva popularitat va ser tan gran que va arribar a aconseguir fins a un 90\% de la quota de mercat en el seu moment més glorios.

Característiques principals:
\begin{itemize}
 \item Interfície gràfica basada en finestres, que simplifica l'ús i redueix la necessitat de comandes de text.
 \item Aplicacions integrades amb funcionalitats com la barra de tasques i el menú d'inici.
 \item Amplia disponibilitat amb programari i facilitat per instal·lar aplicacions.
\end{itemize}

\subsection{Linux}

\label{subsec:Linux}
Linux és un sistema operatiu de codi obert creat per Linus Torvalds el 1991. És conegut per la seva versabilitat.

Característiques principals:
\begin{itemize}
 \item Llibertat i personalització, gràcies al seu model de codi obert.
 \item Gran seguretat i estabilitat.
 \item Gran varietat de  recursos, ideal per a ordinadors antics o tasques especialitzades.
\end{itemize}


\subsection{macOS}

\label{subsec:Mac OS}
macOS és el sistema operatiu exclusiu dels ordinadors Apple, desenvolupat des dels anys 70. Destaca pel seu disseny pulit i la integració amb l'ecosistema de dispositius de la marca.
Característiques principals:
\begin{itemize}
 \item Interfície elegant i intuïtiva
 \item Optimització per al maquinari d'Apple
\end{itemize}

\textbf{Decisió Final: Windows i Linux}

Hem optat per utilitzar Windows i Linux pels següents motius:
\begin{itemize}
\item Linux és ideal per a tasques avançades, terminals i automatitzacions, a més de ser compatible amb el processador Kile.

\item Windows ofereix una major facilitat d'ús i compatibilitat amb programari general, sent perfecte per a tasques diàries.
\end{itemize}

No obstant, en un ordinador rarament hi han dos sistemes operatius convivint, això ens porta ha fer un altre recerca, i n'hem trobat dues:
\begin{enumerate}
 \item Dual boot: Instal·lar amb dues sistemes en particions separades i triar-ne una en l’arrencada.
 \item Màquina virtual: Executar un sistema dins de l’altre.
\end{enumerate}

Finalment, hem escollit la màquina virtual per la seva facilitat d’ús i baixos requisits, encara que el rendiment no sigui tan fluid com en dual boot. Aquesta opció ens permet provar Linux sense afectar la instal·lació principal de Windows.

\subsection{Dos maquines virtual en un sol equip}

\begin{enumerate}
 \item Dual Boot:
      \begin{itemize}
       \item Requereix particionar el disc dur.
       \item Permet un rendiment complet de cada sistema.
       \item Menys flexible per canviar entre sistemes sense reiniciar.
      \end{itemize}

  \item Maquina virtual:
        \begin{itemize}
         \item Menys eficient en recursos (comparteix CPU, RAM i emmagatzematge).
         \item Més fàcil de configurar i gestionar.
         \item Ideal per proves o ús casual sense modificar l’entorn principal.
        \end{itemize}


\end{enumerate}









